\chapter{Einleitung} \label{chpt:Einleitung_Main}

Die Matrixmultiplikation findet in Gebieten Anwendung, die auf massiver Datenverarbeitung und linearer Algebra basieren. Beispiele sind Neuronale Netze und Large Language Models (\gls{LLM}) als eine Anwendung von neuronalen Netzen, sowie die 3D-Computergrafik. Dadurch sind Matrixmultiplikationen ein extrem relevantes Leistungsmerkmal für moderne CPU Architekturen.

In dieser Arbeit soll ein bestehendes RISC-V CPU Design um Anweisungen zur Berechnung der Matrixmultiplikation erweitert werden.

Zunächst wird vertieft dargestellt, weshalb die RISC-V Architektur und die Matrixmultiplikation relevant sind. Dann wird die Aufgabenstellung an sich beschrieben. Es folgt eine Betrachtung des allgemeinen Entwicklungsstandes, in dem die Arbeit ausgeführt wird. Der Stand der Technik wird beschrieben. Das Kapitel wird mit einer Erläuterung zum Aufbau der Arbeit abgeschlossen.




\section{Motivation} \label{sec:Einleitung_Motivation}

In dieser Arbeit soll eine bestehende RISC-V CPU um Anweisungen zur Durchführung von Matrixmultiplikationen erweitert werden.

RISC-V ist eine lizenzfreie Beschreibung eines Prozessors und der Anweisungen, die der Prozessor verarbeiten können muss. Durch diese Lizenzfreiheit ist die RISC-V Architektur sehr interessant, sowohl für den akademischen als auch für den kommerziellen Einsatz.

Die RISC-V Architektur ist die fünfte Iteration des Berkeley RISC Projekts der Universität von Kalifornien. Die Arbeit an RISC-V began im Jahre 2010 und ist inzwischen so ausgereift, dass seit einigen Jahren kommerzielle RISC-V Chips auf dem Markt verfügbar sind.

Die ersten verfügbaren Chips waren aufgrund der fehlenden Lizenzkosten und einfachen Umsetzbarkeit für wenige Cent verfügbar. Im Jahre 2026 gibt es RISC-V Chips, die erweiterte Funktionen bieten und daher sowohl Desktop-Betriebssysteme als auch LLMs ausführen können.

Die Firma Qualcomm beispielsweise, als eines der weltweit führenden Unternehmen auf dem Gebiet der Halbleiter und der Telekommunikation, hat eine RISC-V Erweiterung namens Xqci zu RISC-V und auch zu dem Compiler-Framework LLVM beigesteuert \cite{phoronix_xqci}. Qualcomm führt auch Käufe von Firmen mit RISC-V Expertise durch. Dies deutet daraufhin, dass Qualcomm den Einsatz von RISC-V in Produkten ernsthaft in Erwägung zieht.

Die Matrixmultiplikation ist die Grundlage für perspektivische Projektionen, Rotationen und Bewegung von Objekten in der 3D Computer-Grafik. Dort ist es notwendig 4x4 Matrizen zu verwenden und algebraische Berechnungen so schnell wie möglich auf diesen Matrizen ausführen zu können.

\begin{figure}[h]
\centering
\includegraphics[scale=0.65]{"ViewFrustum.png"}
\caption{View Frustum}
\label{fig:ViewFrustum}
\end{figure}

In Abbildung \ref{fig:ViewFrustum} aus \cite{RTR3} wird auf der linken Seite das Sichtfeld (View Frustum) einer Kamera in einer Szene dargestellt. Durch die Matrix \(P_{p}\) wird dieses Sichtfeld perspektivisch korrekt auf den Einheitswürfel auf der rechten Seite abgebildet. Durch das Umformen des Sichtfelds erscheinen weiter entfernte Objekte z.B. kleiner.

Diese Projektion muss auf jedes Objekt einer Szene angewendet werden. Aus mathematischer Sicht werden also viele Vektoren mit der Projektionsmatrix multipliziert um die transformierten Punkte zu erhalten, die später auf einem Computerbildschirm erscheinen sollen.

Berücksichtigt man die modernen 4K- oder 8K-Auflösungen und die hohen Bildwiederholraten von bis zu 200 Hz, wird deutlich, dass die Matrixmultiplikation extrem schnell und womöglich stark parallelisiert ablaufen muss um eine hohe Bildgenerierungsrate zu erreichen.

Die Projektionsmatrix ist in Abbildung \ref{fig:PerspectiveProjection} aus \cite{RTR3} dargestellt.

\begin{figure}[h]
\centering
\includegraphics[scale=0.7]{"PerspectiveProjection.png"}
\caption{Perspective Projection Matrix}
\label{fig:PerspectiveProjection}
\end{figure}

Der Einsatz von 3D-Computergrafik ist seit mehreren Jahren nicht mehr aus dem Alltag, aber auch aus vielen professionellen Bereichen wegzudenken. Man denke dabei an die Computerspiel- und Film-Industrie sowie Modellierungssoftware für Ingenieursanwendungen aller Art.

Seit Google im Jahr 2017 die Transformer-Architektur erfunden und daraufhin OpenAI ChatGPT 1 programmiert hat, setzt sich die Erfolgsgeschichte von Large Language Models bis heute fort. Generative künstliche Intelligenz wird als revolutionäre Erfindung betrachtet und führt schon jetzt dazu, dass sich Softwareentwickler Gedanken um ihren zukünftigen Wert auf dem Arbeitsmarkt machen. Es ist spätestens im Jahre 2026 klar, dass eine KI bestimmte Aufgaben bedeutend viel schneller erledigt, als es ein Mensch jemals könnte.

Firmen haben diesen Vorteil der KI erkannt und haben vor allem die hohen Lohnkosten menschlicher Arbeit im Blick. Hardware, die \gls{LLM}s trainieren und ausführen kann, wird damit zum relevanten Posten in der Bilanz großer Unternehmen. Anstatt fortlaufende Personalkosten zu bezahlen, möchten Unternehmen KI-fähige Hardware einsetzen, die naiv betrachtet nur einmalige Anschaffungskosten verursacht.

Es soll verdeutlicht werden, welcher Zusammenhang zwischen der Matrixmultiplikation und Neuronalen Netzen bzw. \gls{LLM}s besteht.

Ein LLM basiert auf Neuronalen Netzen. Zwei Neuronen in einem sehr einfachen Netz sind in Abbildung \ref{fig:SimpleNeuralNetzwork} aus \cite{basic-neural-network-matrix-maths-part-1} dargestellt. Dabei sind \(w_{i,j}\) die Gewichte und \(b_{i}\) die Biases, deren Werte durch Training über einer Datenbasis bestimmt werden. Nach dem Training soll das Netz auf Eingaben \(x_{i}\) mit sinnvollen Ausgabewerten \(y_{i}\) reagieren.

Zusammen werden Gewichte und Biases als Parameter bezeichnet. Das kleinste frei verfügbare LLM verwendet 10 Millionen Parameter, große LLM verwenden Billionen von Parametern.

\begin{figure}[h]
\centering
\includegraphics[scale=0.9]{"LLM_1.png"}
\caption{Einfaches Neuronales Netz}
\label{fig:SimpleNeuralNetzwork}
\end{figure}

Ein Neuronales Netz lässt sich als Matrix beschreiben. Eingaben an ein Neuronales Netz erfolgen in Form von Werten \(x_{i,j}\). Werden mehrere Eingaben \(x_{i}\) gleichzeitig an das neuronale Netz gegeben, wird aus den einzelnen Eingabevektoren eine Matrix. Dieser Umstand ist in Abbildung \ref{fig:InputBatching} aus \cite{basic-neural-network-matrix-maths-part-1} dargestellt.

\begin{figure}[h]
\centering
\includegraphics[scale=0.55]{"LLM_2.png"}
\caption{Mehrere kombinierte Eingaben}
\label{fig:InputBatching}
\end{figure}

Die rasante Verbreitung von LLMs hat Auswirkungen auf die Hardwarehersteller und den Handel der Hardware. Die Firma Micron hat beispielsweise einen Rückzug aus dem Endkundengeschäft vollzogen und verkauft RAM-Bausteine nur noch an Firmenkunden, die den RAM-Speicher dann bevorzugt für die Berechnung in Rechenzentren einsetzen. Es liegt nahe, dass LLMs in den Rechenzentren trainiert und betrieben werden. Die RAM-Preise für Endkunden sind aus diesem Grund innerhalb der Jahre 2025 und 2026 um ein Vielfaches angestiegen.

Gleiches gilt für Grafikbeschleuniger der Firma NVIDIA. Auch hier konkurriert der Endkunde mit dem Firmenkunden. NVIDIA kann die Preise beinahe beliebig festlegen. Das Top-Produkt für Endkunden der Firma NVIDIA ist eine 5090 RTX Grafikkarte. Sie ist im Januar 2026 für etwa 3600 Euro erhältlich und kostet damit mehr als der gesamte Rest eines Computers. Eine Grafikkarte ist in der Lage Matrixmultiplikationen in großer Parallelität und extrem schnell zu berechnen. Grafikkarten werden auch zum Training von LLMs verwendet.

Getrieben durch den Erfolg von LLMs kann sich eine Spezifikation für CPUs, wie z.B. RISC-V, nicht mehr durchsetzen, wenn sie keine Anweisungen für den effizienten Umgang mit Vektoren und Matrizen besitzt. Die Endkunden werden Produkte, die nicht in der Lage sind, Matrizen hardwarebeschleunigt zu berechnen, nicht mehr kaufen, z.B. aus Gründen der Computergrafik. Das gleiche gilt für Firmenkunden, die den Betrieb eines LLM planen. Ohne beschleunigte Matrixberechnung kann sich eine CPU nicht mehr auf dem Markt behaupten.

Die CPUs der Firmen AMD und Intel besitzen bereits die Erweiterungen AVX2 (Advanced Vector Extensions 2), AMX (Advanced Matrix Extensions) und FMA (Fused Multiply Add) zum Rechnen mit Vektoren und Matrizen.
% https://de.wikipedia.org/wiki/FMA_x86

All diese Punkte zusammengenommen ergeben eine große Motivation für den Einsatz von RISC-V und den Einsatz von Matrixmultiplikation innerhalb einer RISC-V-CPU.




\section{Aufgabenstellung} \label{sec:Einleitung_Aufgabenstellung}

Die NEORV32-CPU wird um Anweisungen erweitert, die die hardwarebeschleunigte Matrixmultiplikation ermöglichen. Die NEORV32 CPU liegt in Form einer textuellen Beschreibung vor. Diese Beschreibung wird ergänzt, durch ein Werkzeug übersetzt und kann dann auf einem FPGA ausgeführt werden.

Ein Field Programmable Gate Array (FPGA) ist ein Chip, der eine große Menge an logischen Gattern bereitstellt, die zunächst nicht miteinander verbunden sind, also zunächst keine Funktion ausführen. Das Innere des FPGA muss programmiert werden. Damit kann ein FPGA unterschiedliche Hardware abbilden, je nachdem was die textuelle Beschreibung vorgibt.

Die neu hinzugefügten Anweisungen orientieren sich von der Kleinteiligkeit her an den Anweisungen, wie sie bereits für die \gls{RVV}-Vektor-Extension in RISC-V spezifiziert worden sind. Die Vektor-Extension wird im Abschnitt \ref{sec:Stand_der_Technik_Main} beschrieben. Im Speziellen wird das Konzept des Strip-Mining ermöglicht. Dabei handelt es sich um die Aufteilung des Problems in Teile, die der CPU, aufgrund von begrenzten Ressourcen, schrittweise zugeführt werden, bis die Aufgabe als Ganzes gelöst ist.

Die Erweiterung soll durch Entwickler mit geringer Erfahrung auf dem Feld der Hardware"=Beschreibung durchgeführt werden können. Die Matrizenberechnungen erfolgen daher lediglich auf Matrizen, deren Dimensionen ein Vielfaches der Zahl vier sind (also z.B. 4x4, 8x8, 16x16) und es werden nur ganzzahlige Multiplikationen ohne Vorzeichen durchgeführt. Ohne Beschränkung der Allgemeinheit lässt sich eine Matrix an den Rändern mit Nullen auffüllen um auf ein Vielfaches der Dimension vier zu kommen, ohne das Ergebnis der Multiplikation zu verändern.

Da es viele Ansätze gibt Matrizen zu multiplizieren, wird eine Einschränkung vorgegeben. Die Multiplikation erfolgt mit dem Outer Product, da dieser Ansatz auch in dem maßgebenden Dokument zu dieser Arbeit beschrieben ist \cite{design_VME}. Der deutsche Begriff für das Outer Product ist das dyadische Produkt.




\section{Stand der Technik} \label{sec:Stand_der_Technik_Main}

%\chapter{Stand der Technik} \label{chpt:Stand_der_Technik_Main}

Ein RISC-V Komitee hat im Jahre 2015 mit der ersten Version der RISC-V Vector Extension (\gls{RVV}) begonnen. Im Jahre 2021 wurde das Release 1.0 ratifiziert und ist jetzt in einen finalen Zustand (Frozen) überführt worden. RISC-V wird durch die (\gls{RVV}) um Register für Vektoren und explizite Anweisungen zur Manipulation dieser Vektoren erweitert.

%Diese führte dazu, dass bereits Versuche unternommen worden sind, kleine lokale Large-Language-Modelle auf RISC-V CPUs auszuführen. Dabei zeigt sich, dass sich kleine LLMs tatsächlich mit ? Token pro Sekunde unterhalten können.

In der RISC-V Vector Extension Spezifikation \cite{rvv_specification} wird das Konzept des Strip-Mining beschrieben und hervorgehoben. Die Software lädt den Cache mit einem Teil einer oder mehrerer Vektoren voll und beginnt den Cache dann abzuarbeiten, in dem der Teil der Vektoren dann wiederum teilweise in den CPU geladen, verarbeitet und danach wieder aus dem CPU zurück in den Speicher übertragen wird. Das Strip-Mining wiederholt diese Abläufe, bis die Aufgabe erledigt ist, also alle Element der Vektoren bearbeitet sind. Das Beispiel in Kapitel 6.4. namens ``Example of stripmining and changes to SEW'' aus \cite{rvv_specification} verdeutlicht dieses Vorgehen.

%Diese zwei Abläufe (Cache-Frame laden und Register laden) werden dann solange wiederholt, bis der gesamte Vektor abgearbeitet ist. Ein Hinweis auf dieses Strip-Mining Vorgehen findet sich in der RISC-V Vector Extension Spezifikation. zum Strip-Mining.

Die Alternative zum Hardware unterstützen Strip-Mining wäre es, jedes einzelne Vektorelement einzeln in ein Register der CPU zu laden und dann mit einem Element des zweiten Vektors zu verrechnen. Dabei wird für jedes Element ein Laden und Speicher aus dem Speicher notwendig und man bezahlt mit einem hohen Verlust an Geschwindigkeit für die im Gegenzug triviale Implementierung.

Der Vorteil der \gls{RVV} Register ist es, dass diese Vektorregister mehrere Elemente parallel auf eine Instruktion hin zusammenrechnen können. Der CPU verfährt dann in einem echten \gls{SIMD}-Betriebsmodus (Single Instruction Multiple Data) und spart ein vielfaches an Taktzyklen im Vergleich zur sequentiellen Ausführung.

%Es sind auch komplett neuartige Designs denkbar. Wenn der Cache-Speicher über mehrere Kanäle an den CPU angebunden ist, könnte der CPU die Vektorregister auch parallel über alle Kanäle befüllen, mehrer Register parallel verrechnen und die Ergebnisse parallel wieder in den Cache zurückliefern.

% https://github.com/flame/blis/blob/master/docs/KernelsHowTo.md

Das Prinzip des Strip-Mining wird auch in mathematischen Bibliotheken verwendet. Eine bekannte Bibliothek ist blis \cite{BLIS1}. blis verwendet den Begriff Kernel oder Microkernel. Ein Kernel ist Code, der auf einem Cache-Frame ausgeführt wird, also den Code darstellt, den die CPU sehr schnell ausführen kann. Optimierte Anweisung für verschiedene CPU Architekturen werden hauptsächlich in dem Code verwendet, der dem Kernel zuzurechnen ist um Zeit zu sparen.

Um die Berechnungen im Kernel herum gibt es weiteren Code, der über Schleifen nacheinander alle Daten in den Kernel lädt, bis die Berechnung vollständig abgeschlossen ist.

Es gibt sehr starke Parallelen zwischen dem Strip-Mining und dem Kernel-Ansatz der blis Bibliothek. Sollte blis für RISC-V portiert werden, ist es sinnvoll, dass RISC-V das Prinzip des Strip-Mining ermöglicht. Dies wurde durch die RISC-V Vector Extension für Vektoren bereits erreicht.

%blis dokumentiert die Matrixmultiplikation unter dem Stichwort GEMM. GEMM steht für General Matrix Multiply. Der gemm kernel wird für die GEMM Algorithmen verwendet.

Für die RISC-V Matrixberechnung gibt es zum jetzigen Zeitpunkt keine ratifizierte Extension. Innerhalb der RISC-V Arbeitsgruppen werden unterschiedliche Ansätze diskutiert. Die Ansätze werden im Kapitel zu den theoretischen Grundlagen beschrieben.

Interessanterweise gibt es eine chinesische Firma namens Stream Computing die sich zur Mission gesetzt hat, high-end RISC-V Chips zum Einsatz in Rechenzentren zu entwickeln. Bereits im Jahre 2022 konnte Stream Computing ein eigenes Arbeitsergebnis vorweisen \footnote{https://riscv.org/blog/stream-computing-risc-v-matrix-extension-open-source-project-upgrades-to-version-0-5-supporting-vectormatrix-implementation/}. Dieses Arbeitsergebnis wurde von Stream Computing selbst ``RISC-V Matrix ISA Specification'' in der Version v0.1 genannt. Die Version v0.5 wurde im Oktober 2024 erstellt. Es handelt sich aber trotz des Namens nicht um die offizielle Version, die das RISC-V Gremium veröffentlichen wird! Stream Computing hat ebenfalls die notwendigen Anpassungen zum LLVM Compiler-Framework beigetragen!

Wie tragbar die Ergebnisse sind lässt sich schwer einschätzen ohne eine eingehende Evaluierung des Quellcodes durchgeführt zu haben. Es ist jedoch interessant zu sehen, welche Innovationskraft und Arbeitsgeschwindigkeit die chinesischen Industrie hat. Während sich das RISC-V Gremium wieder auf eine mehrjährige Reise begibt um eine Matrix-Extension zu erzeugen, wie es bereits bei der Vektor-Extension erfolgt ist, hat eine chinesische Firma vermeintlich bereits ein Design und eine Toolchain zu diesem Design erstellt.

% TODO Verweis: https://riscv.org/blog/stream-computing-risc-v-matrix-extension-open-source-project-upgrades-to-version-0-5-supporting-vectormatrix-implementation/

Innerhalb des NEORV32 CPU Umfeldes gibt es eine Arbeit von Qizal Arsalan \cite{QArsalan_NEORV32_MATRIX_CFU} bei der der NEORV32 CPU um Matrixberechnungen erweitert wird. Diese Bachelorarbeit unterscheidet sich allerdings hinsichtlich der Herangehensweise stark von Qizal Arsalans Arbeit. Der NEORV32 CPU besitzt im Design vorgesehene Erweiterungspunkte namens CFS und CFU. Das CFS (Custom Functions Subsystem) erlaubt es Hardware-Beschleuniger in den NEORV32 CPU einzubringen, die über Speicheradressen (Memory Mapping) aus einer Softwareanwendung mit Daten versorgt werden können. CFU (Custom Function Unit) erlaubt es einzelne Instruktionen zu ergänzen um den CPU um kleinschrittige Operationen zu erweitern.

Qizal Arsalan untersucht den Ansatz der Matrixmultiplikation mit CFS und CFU. In dieser Bachelorarbeit werden Matrix-Instruktionen in den NEORV32 Chip eingebracht, die nicht über das CFS oder CFU Feature eingebunden werden sondern über die Execution-Engine, also der State-Machine, des NEORV32 Chips selbst. Der Grund ist, dass die Speicheraddressabhängigkeit von CFS und CFU nicht konform mit einer Matrix-Extension-Spezifikation des RISC-V Gremiums sein können und nicht sein werden.

Bei CFS und CFU handelt es sich um Systeme, die nicht Teil der RISC-V Spezifikation sind. Diese Bachelorarbeit versucht die Ergebnisse einer Matrix-Extension im RISC-V Umfeld zu antizipieren und kann daher CFS und CFU nicht verwenden.
\chapter{Stand der Technik} \label{chpt:Stand_der_Technik_Main}

Ein RISC-V Komitee hat im Jahre 2015 mit der ersten Version der RISC-V Vector Extension (\gls{RVV}) begonnen. Im Jahre 2021 wurde das Release 1.0 ratifiziert und ist jetzt in einen finalen Zustand (Frozen) überführt worden. RISC-V wird durch die (\gls{RVV}) um Register für Vektoren und explizite Anweisungen zur Manipulation dieser Vektoren erweitert.

%Diese führte dazu, dass bereits Versuche unternommen worden sind, kleine lokale Large-Language-Modelle auf RISC-V CPUs auszuführen. Dabei zeigt sich, dass sich kleine LLMs tatsächlich mit ? Token pro Sekunde unterhalten können.

In der RISC-V Vector Extension Spezifikation \cite{rvv_specification} wird das Konzept des Strip-Mining beschrieben und hervorgehoben. Die Software lädt den Cache mit einem Teil einer oder mehrerer Vektoren voll und beginnt den Cache dann abzuarbeiten, in dem der Teil der Vektoren dann wiederum teilweise in den CPU geladen, verarbeitet und danach wieder aus dem CPU zurück in den Speicher übertragen wird. Das Strip-Mining wiederholt diese Abläufe, bis die Aufgabe erledigt ist, also alle Element der Vektoren bearbeitet sind. Das Beispiel in Kapitel 6.4. namens ``Example of stripmining and changes to SEW'' aus \cite{rvv_specification} verdeutlicht dieses Vorgehen.

%Diese zwei Abläufe (Cache-Frame laden und Register laden) werden dann solange wiederholt, bis der gesamte Vektor abgearbeitet ist. Ein Hinweis auf dieses Strip-Mining Vorgehen findet sich in der RISC-V Vector Extension Spezifikation. zum Strip-Mining.

Die Alternative zum Hardware unterstützen Strip-Mining wäre es, jedes einzelne Vektorelement einzeln in ein Register der CPU zu laden und dann mit einem Element des zweiten Vektors zu verrechnen. Dabei wird für jedes Element ein Laden und Speicher aus dem Speicher notwendig und man bezahlt mit einem hohen Verlust an Geschwindigkeit für die im Gegenzug triviale Implementierung.

Der Vorteil der \gls{RVV} Register ist es, dass diese Vektorregister mehrere Elemente parallel auf eine Instruktion hin zusammenrechnen können. Der CPU verfährt dann in einem echten \gls{SIMD}-Betriebsmodus (Single Instruction Multiple Data) und spart ein vielfaches an Taktzyklen im Vergleich zur sequentiellen Ausführung.

%Es sind auch komplett neuartige Designs denkbar. Wenn der Cache-Speicher über mehrere Kanäle an den CPU angebunden ist, könnte der CPU die Vektorregister auch parallel über alle Kanäle befüllen, mehrer Register parallel verrechnen und die Ergebnisse parallel wieder in den Cache zurückliefern.

% https://github.com/flame/blis/blob/master/docs/KernelsHowTo.md

Das Prinzip des Strip-Mining wird auch in mathematischen Bibliotheken verwendet. Eine bekannte Bibliothek ist blis \cite{BLIS1}. blis verwendet den Begriff Kernel oder Microkernel. Ein Kernel ist Code, der auf einem Cache-Frame ausgeführt wird, also den Code darstellt, den die CPU sehr schnell ausführen kann. Optimierte Anweisung für verschiedene CPU Architekturen werden hauptsächlich in dem Code verwendet, der dem Kernel zuzurechnen ist um Zeit zu sparen.

Um die Berechnungen im Kernel herum gibt es weiteren Code, der über Schleifen nacheinander alle Daten in den Kernel lädt, bis die Berechnung vollständig abgeschlossen ist.

Es gibt sehr starke Parallelen zwischen dem Strip-Mining und dem Kernel-Ansatz der blis Bibliothek. Sollte blis für RISC-V portiert werden, ist es sinnvoll, dass RISC-V das Prinzip des Strip-Mining ermöglicht. Dies wurde durch die RISC-V Vector Extension für Vektoren bereits erreicht.

%blis dokumentiert die Matrixmultiplikation unter dem Stichwort GEMM. GEMM steht für General Matrix Multiply. Der gemm kernel wird für die GEMM Algorithmen verwendet.

Für die RISC-V Matrixberechnung gibt es zum jetzigen Zeitpunkt keine ratifizierte Extension. Innerhalb der RISC-V Arbeitsgruppen werden unterschiedliche Ansätze diskutiert. Die Ansätze werden im Kapitel zu den theoretischen Grundlagen beschrieben.

Interessanterweise gibt es eine chinesische Firma namens Stream Computing die sich zur Mission gesetzt hat, high-end RISC-V Chips zum Einsatz in Rechenzentren zu entwickeln. Bereits im Jahre 2022 konnte Stream Computing ein eigenes Arbeitsergebnis vorweisen \footnote{https://riscv.org/blog/stream-computing-risc-v-matrix-extension-open-source-project-upgrades-to-version-0-5-supporting-vectormatrix-implementation/}. Dieses Arbeitsergebnis wurde von Stream Computing selbst ``RISC-V Matrix ISA Specification'' in der Version v0.1 genannt. Die Version v0.5 wurde im Oktober 2024 erstellt. Es handelt sich aber trotz des Namens nicht um die offizielle Version, die das RISC-V Gremium veröffentlichen wird! Stream Computing hat ebenfalls die notwendigen Anpassungen zum LLVM Compiler-Framework beigetragen!

Wie tragbar die Ergebnisse sind lässt sich schwer einschätzen ohne eine eingehende Evaluierung des Quellcodes durchgeführt zu haben. Es ist jedoch interessant zu sehen, welche Innovationskraft und Arbeitsgeschwindigkeit die chinesischen Industrie hat. Während sich das RISC-V Gremium wieder auf eine mehrjährige Reise begibt um eine Matrix-Extension zu erzeugen, wie es bereits bei der Vektor-Extension erfolgt ist, hat eine chinesische Firma vermeintlich bereits ein Design und eine Toolchain zu diesem Design erstellt.

% TODO Verweis: https://riscv.org/blog/stream-computing-risc-v-matrix-extension-open-source-project-upgrades-to-version-0-5-supporting-vectormatrix-implementation/

Innerhalb des NEORV32 CPU Umfeldes gibt es eine Arbeit von Qizal Arsalan \cite{QArsalan_NEORV32_MATRIX_CFU} bei der der NEORV32 CPU um Matrixberechnungen erweitert wird. Diese Bachelorarbeit unterscheidet sich allerdings hinsichtlich der Herangehensweise stark von Qizal Arsalans Arbeit. Der NEORV32 CPU besitzt im Design vorgesehene Erweiterungspunkte namens CFS und CFU. Das CFS (Custom Functions Subsystem) erlaubt es Hardware-Beschleuniger in den NEORV32 CPU einzubringen, die über Speicheradressen (Memory Mapping) aus einer Softwareanwendung mit Daten versorgt werden können. CFU (Custom Function Unit) erlaubt es einzelne Instruktionen zu ergänzen um den CPU um kleinschrittige Operationen zu erweitern.

Qizal Arsalan untersucht den Ansatz der Matrixmultiplikation mit CFS und CFU. In dieser Bachelorarbeit werden Matrix-Instruktionen in den NEORV32 Chip eingebracht, die nicht über das CFS oder CFU Feature eingebunden werden sondern über die Execution-Engine, also der State-Machine, des NEORV32 Chips selbst. Der Grund ist, dass die Speicheraddressabhängigkeit von CFS und CFU nicht konform mit einer Matrix-Extension-Spezifikation des RISC-V Gremiums sein können und nicht sein werden.

Bei CFS und CFU handelt es sich um Systeme, die nicht Teil der RISC-V Spezifikation sind. Diese Bachelorarbeit versucht die Ergebnisse einer Matrix-Extension im RISC-V Umfeld zu antizipieren und kann daher CFS und CFU nicht verwenden.




\section{Struktur der Arbeit} \label{sec:Einleitung_Struktur_der_Arbeit}

Die Matrixmultiplikation wird im Kapitel ``Theoretische Grundlagen'' beschrieben. Es wird verdeutlicht, dass sich die Matrixmultiplikation durch Aufteilung der Matrizen in Blöcke aufteilen und damit schrittweise abarbeiten lässt. Damit wird das Strip-Mining ermöglicht.

%Nachdem die Matrixmultiplikation beschrieben ist, wird ihre Relevanz in den Anwendungsbegieten der Neuronalen Netze und der 3D-Computergrafik dargestellt.

Danach werden die momentan ablaufenden Anstrengungen des RISC-V Gremiums zur Erstellung einer Extension zur Matrixmultiplikation beschrieben. Dabei gibt es drei Ansätze. Der in dieser Bachelorarbeit untersuchte Ansatz wird zusammen mit den anderen Ansätzen vorgestellt. Dabei wird der Strip-Mining-Ansatz erläutert und es wird eine Softwareanwendung besprochen, die diesen Ansatz als Computerprogramm implementiert.

Es folgt eine Beschreibung des Vorgehens bei der Erstellung dieser Bachelor-Arbeit.

Nachdem der Ansatz beschrieben ist, wird die NEORV32 CPU als das Ziel der Implementierung beschrieben. Das Design wird herausgearbeitet und es wird diskutiert, welche Änderungen an welchen Stellen vorgenommen werden, um die CPU um Matrixbefehle zu erweitern.

Danach wird der VHDL Quellcode der NEORV32 CPU erweitert und auf einen FPGA programmiert. Auf dem FPGA wird eine Software-Anwendung ausgeführt, die die Matrix-Instruktionen verwendet um eine Matrixmultiplikation durchzuführen.

Schließlich wird das Ergebnis mit einer normalen Matrixberechnung verglichen.

Die Arbeit endet mit einer Zusammenfassung und einem Ausblick auf mögliche Verbesserungen.

Der Anhang enthält die Design-Dateien aus dem NEORV32 Projekt, die im Zuge dieser Arbeit geändert worden sind.